\chapter*{Bibliographie générale}
\addcontentsline{toc}{chapter}{Bibliographie générale}
\markboth{Bibliographie générale}{Bibliographie générale}
Les références ci-dessous sont proposées comme ouvrages de travail et ressources d'approfondissement pour l'ensemble du cours. Les exercices du présent polycopié sont reformulés à partir de ces sources.
\begin{itemize}[leftmargin=0pt,label={},itemsep=9pt]
\item R. Cori et D. Lascar, \emph{Logique mathématique, tome 1 : Calcul propositionnel, algèbre de Boole, calcul des prédicats}, Dunod, Paris, 2003.
\item S. Epp, \emph{Discrete Mathematics with Applications}, 4th ed., Brooks/Cole, 2011.
\item H. B. Enderton, \emph{Elements of Set Theory}, Academic Press, New York, 1977.
\item P. R. Halmos, \emph{Naive Set Theory}, D. Van Nostrand, Princeton, 1960.
\item R. Hammack, \emph{Book of Proof}, 3rd ed., 2018 (version révisée en 2025). Version PDF libre : \url{https://richardhammack.github.io/BookOfProof/}.
\item J.-L. Krivine, \emph{Théorie des ensembles}, Cassini, Paris, 1998.
\item E. Lehman, F. T. Leighton et A. R. Meyer, \emph{Mathematics for Computer Science}, MIT OpenCourseWare, cours en ligne : \href{https://ocw.mit.edu/courses/6-1200j-mathematics-for-computer-science-spring-2024/}{MIT OpenCourseWare}.
\item K. H. Rosen, \emph{Discrete Mathematics and Its Applications}, 8th ed., McGraw-Hill, 2019.
\item D. J. Velleman, \emph{How to Prove It: A Structured Approach}, 3rd ed., Cambridge University Press, 2019.
\item R. Chemlal, \emph{Cours de logique mathématique}, polycopié de Licence 2, Université de Béjaïa, version institutionnelle : \href{https://elearning.univ-bejaia.dz/pluginfile.php/503483/mod_resource/content/0/Cours_CHEMLAL%20Rezki_Logique%20Math%C3%A9matique.pdf}{Université de Béjaïa}.
\end{itemize}
